\clearpage\fancyhead[L]{\small PALLAS.}\hypertarget{native-388-lower}{}\label{native-388-lower}
\begin{center}\textbf{PALLAS.}\par\vspace{4pt}\hrule\vspace{4pt}\fontsize{10}{12}\selectfont Astronomisches Jahrbuch für 1818. pp. 167\;..\;173. Berlin 1815.\par\vspace{4pt}\hrule\vspace{4pt}\end{center}

\begin{flushright}
Göttingen, 1 February 1815.
\end{flushright}

\par\indent The calculations for \emph{Pallas}, which also appear in the last yearbook, have once again been borne out remarkably well. I myself observed the planet only once with the ring micrometer

\begin{center}
\fontsize{9}{11}\selectfont\setlength{\tabcolsep}{3pt}\renewcommand{\arraystretch}{1.15}\begin{tabular}{l|ccc|ccc|ccc}
1814 & \multicolumn{3}{c|}{Mean time} & \multicolumn{3}{c|}{Right ascension} & \multicolumn{3}{c}{Declination} \\
\hline
Sept.\ 16 & $12^{\mathrm u}$ & $1^{\mathrm m}$ & $1^{\mathrm s}$ & $146^\circ$ & $16'$ & $23''4$ & $11^\circ$ & $4'$ & $20''9$ S. \\
\end{tabular}
\end{center}

\noindent but a valuable series of meridian observations by Messrs. von \textsc{Lindenau}, \textsc{Bessel} and \textsc{Schumacher} enabled me to calculate the opposition accurately. For details I refer to the Gött.\ gel.\ Anz., no.\ 199 [p.\ 379 of this volume]. I give only the final result here:

\begin{center}25 October 1814. $12^{\mathrm u}\;33^{\mathrm m}\;22^{\mathrm s}$ Mean time in Göttingen\end{center}
\begin{center}\fontsize{10}{12}\selectfont\setlength{\tabcolsep}{3pt}\renewcommand{\arraystretch}{1.15}\begin{tabular}{lcr}
True longitude & $.\;.\;.\;.\;.\;.\;.\;.\;.\;.\;.\;.\;.$ & $31^\circ\;58'\;11''3$\\
Geocentric latitude, south & $.\;.\;.\;.\;.\;.\;.\;.$ & $37\;20\;53.2$\\
\end{tabular}\end{center}

\par\indent My theory accounts for all 10 oppositions so far to within a few seconds.

\Needspace{12\baselineskip}\par\indent Mr. \textsc{Enke} has undertaken the valuable work of making the calculations for the next apparition in advance. I give his result for opposition XI here:

\begin{center}9 April 1816. $2^{\mathrm u}\;14^{\mathrm m}\;24^{\mathrm s}$ Mean time in Göttingen\end{center}
\begin{center}\fontsize{10}{12}\selectfont\setlength{\tabcolsep}{3pt}\renewcommand{\arraystretch}{1.15}\begin{tabular}{lcr}
True longitude & $.\;.\;.\;.\;.\;.\;.\;.\;.\;.\;.\;.$ & $199^\circ\;34'\;40''5$\\
Geocentric latitude & $.\;.\;.\;.\;.\;.\;.\;.\;.$ & $28\;6\;16.3\;\text{North}$\\
\end{tabular}\end{center}

\Needspace{8\baselineskip}\par\indent At the next opposition \emph{Pallas} will be brighter than it has been for several years. Here is a comparison of the light intensities

\begin{center}\fontsize{9}{11}\selectfont
1812\ .\ 0.01666,\quad 1813\ .\ 0.01475,\quad 1814\ .\ 0.05476,\quad 1816\ .\ 0.05997
\end{center}

\par\indent I have the pleasure of enclosing a copy of the ephemeris calculated by the same skilled young astronomer.

\par\indent The weather was unusually favorable for observing the winter solstice this year. I spared no effort to make these observations as accurate as possible. One of these measures \SourceLeaf{389}was a separate screen that let light fall only on the objective and kept out all other rays. Since I had yet to fit this device when I observed earlier solstices, I was unwilling to place complete confidence in those observations. Even so, I cannot really say that the observations before the screen was fitted agreed less well. The result for the obliquity of the ecliptic now agrees exactly with Mr. \textsc{von Zach}’s. The winter solstice gives a value $10''$ lower than the summer solstice. I would prefer, however, to express this phenomenon as follows. If we take the mean value for the obliquity of the ecliptic, observations of the Sun always give an elevation of the pole that is on average $5''$ lower than that given by circumpolar stars. In my experience this holds in \emph{every} season. This strange phenomenon is still an astronomical puzzle. It would be solved if there were something in the instrument that justified assuming that it gives all zenith distances $2''5$ too small. But I can find absolutely nothing that would always produce an error in the same direction. I especially want to point out that I aligned the telescope’s line of sight parallel to the plane of the instrument with the utmost accuracy, using a method of my own. I need hardly mention that I also took care to make the plane exactly vertical and to leave nothing out of the calculation. For example, the terms depending on the fourth power were included every time. Neglecting any of these precautions does indeed make the zenith distance too small. My observations of the last solstice were as follows:

\begin{center}
\fontsize{10}{12}\selectfont\setlength{\tabcolsep}{5pt}\renewcommand{\arraystretch}{1.15}\begin{tabular}{l l ccc l}
1814 Dec.\ 17 & T.\ Z.\ D.\ & $74^\circ$ & $53'$ & $48''53$ & 12 obs. \\
19 & & 74 & 57 & 35.13 & 14 \\
21 & & 74 & 59 & 24.02 & 14 \\
22 & & 74 & 59 & 38.63 & 14 \\
23 & & 74 & 59 & 21.10 & 14 \\
24 & & 74 & 58 & 36.74 & 14 \\
25 & & 74 & 57 & 22.29 & 14 \\
29 & & 74 & 57 & 51.39 & 12 \\
\end{tabular}
\end{center}

\par\indent On each day the observations were divided into 2 or 3 series, each of which usually differed from the mean by scarcely one second and never by two.

\par\indent I have yet to find time to calculate these observations with full precision. A preliminary result, however, is a zenith distance of $74^\circ\,59'\,36''6$ at the solstice. With the elevation of the pole $51^\circ\,31'\,55''0$, this gives an obliquity $5''$ smaller than the value used in our best tables. The elevation of the pole is based on more than 200 observations of Polaris at its upper and lower culminations. The latter gave $51^\circ\,31'\,55''5$, the former $51^\circ\,31'\,54''5$, when I used Mr. \textsc{Bessel}’s tables.

\Needspace{17\baselineskip}\begin{flushright}
Göttingen, 24 April
\end{flushright}

\par\indent You will doubtless have found and observed \textsc{Olbers}’s comet yourself long before now. I did find it on 13 March, as soon as I received the first news of its discovery. But the extremely unfavorable weather prevented me from making an actual observation until the 20th. My observations in March are as follows:

\begin{center}\fontsize{9}{11}\selectfont\setlength{\tabcolsep}{3pt}\renewcommand{\arraystretch}{1.15}\begin{tabular}{l|ccc|ccc|ccc}
1815 & \multicolumn{3}{c|}{Mean time} & \multicolumn{3}{c|}{\shortstack{Apparent right\\ascension}} & \multicolumn{3}{c}{\shortstack{Apparent decl.\\North}} \\

\hline
March 20 & $10^{\mathrm u}$ & $33^{\mathrm m}$ & $6^{\mathrm s}$ & $54^\circ$ & $7'$ & $1''$ & $39^\circ$ & $7'$ & $47''$ \\
21 & 10 & 11 & 37 & 54 & 33 & 21 & 39 & 36 & 57 \\
21 & 10 & 57 & 0 & 54 & 34 & 21 & 39 & 38 & 8 \\
25 & 9 & 46 & 38 & 56 & 28 & 50 & 41 & 38 & 5 \\
30 & 9 & 50 & 57 & 59 & 14 & 3 & 44 & 10 & 27 \\
\end{tabular}\end{center}

\Needspace{20\baselineskip}\SourceLeaf{390}\par\indent The second observation on the 21st is by Mr. \textsc{Enke}, who also reduced all the comparison stars from the \emph{Histoire Céleste}. My April observations have yet to be reduced. --- On 31 March I calculated the following preliminary parabolic elements from the observations above. Even now they should differ little.

\begin{center}\fontsize{9}{11}\selectfont $\text{Perihelion}\left\{\begin{array}{lcr}
\text{Time, 24 April 1815} & .\;.\;.\;. & 16^{\mathrm u}\;37^{\mathrm m}\;34^{\mathrm s}\;\text{Göttingen time}\\
\text{Longitude} & .\;.\;.\;.\;.\;.\;.\;.\;. & 146^\circ\;7\;2\\
\text{Distance} & .\;.\;.\;.\;.\;.\;.\;. & 1.24738\\
\end{array}\right.$\end{center}
\begin{center}\fontsize{10}{12}\selectfont\setlength{\tabcolsep}{3pt}\renewcommand{\arraystretch}{1.15}\begin{tabular}{lcr}
Ascending node & $.\;.\;.\;.\;.\;.\;.\;.$ & $82\;43\;4$\\
Inclination of orbit & $.\;.\;.\;.\;.\;.\;.\;.\;.$ & $45\;8\;55$\\
\end{tabular}\end{center}
\noindent Prograde motion.

\par\indent I hope this comet can still be observed until July and that its orbit can then be calculated as an ellipse with considerable confidence.

\par\indent Juno was only 9th to 10th magnitude around the time of its opposition. I observed it well twice, but have yet to reduce the observations.

\par\indent Little use has so far been made of our excellent heliometer for want of a fixed mounting. The day before yesterday, the parallactic mount that had been on order for nearly a year finally arrived. It was very finely and ingeniously made by \textsc{Frauenhofer}.
